LLVManim is implemented as a three-layer pipeline: \textbf{Ingestion}, \textbf{Transformation}, and \textbf{Presentation}. Layer boundaries are enforced at build time by \texttt{import-linter}, ensuring that no module in a lower layer imports from a higher one. The entire codebase is written in Python~3.12 and managed with \texttt{uv}.

\begin{figure}[htbp]
\centerline{\includegraphics[width=\linewidth]{ModelDiagram}}
\caption{Three-layer architecture of LLVManim. The Ingestion layer collects LLVM IR and metadata, the Transformation layer maps it to an intermediate scene graph, and the Presentation layer renders animations or exports artifacts.}
\label{fig:model_diagram}
\end{figure}

\subsubsection{Ingestion Layer}
The Ingestion layer extracts program structure from LLVM IR. It uses the \texttt{llvmlite}\cite{llvmlite} Python binding to parse \texttt{.ll} files and also supports importing pre-computed data from JSON.

\begin{itemize}
    \item \textbf{\texttt{llvm\_events.py}:} Parses an LLVM module and emits a \texttt{ProgramEventStream}---a typed sequence of \texttt{IREvent} records. Each event carries the function name, block name, opcode, full instruction text, a classified \texttt{EventKind} (one of \texttt{alloca}, \texttt{load}, \texttt{store}, \texttt{binop}, \texttt{compare}, \texttt{call}, \texttt{ret}, \texttt{br}, or \texttt{other}), its index within the function, debug line number, and operand list. Typed CFG edges are extracted directly from \texttt{llvmlite} terminator operands (\texttt{br}, \texttt{switch}, \texttt{invoke}, \texttt{indirectbr}, \texttt{callbr}) with true/false labels on conditional branches.

    \item \textbf{\texttt{cfg\_edge\_io.py}:} Serializes and deserializes \texttt{CFGEdge} lists to JSON, enabling round-trip import/export of control-flow edges via the \texttt{--import-cfg-edges} and \texttt{--export-cfg-edges} CLI flags.

    \item \textbf{\texttt{analysis\_metadata\_io.py}:} Loads and saves dominator-tree and loop-analysis metadata as JSON (\texttt{--import-analysis-metadata} / \texttt{--export-analysis-metadata}).

    \item \textbf{\texttt{trace\_io.py}:} Reads and writes runtime path traces as JSON (\texttt{--import-trace} / \texttt{--export-trace}), used for overlay visualization in CFG traversal animations.

    \item \textbf{\texttt{dot\_layout.py}:} Parses a Graphviz DOT file produced by \texttt{opt -passes=dot-cfg} and computes a \texttt{DotLayout} with node positions and edge routes for the CFG animation scene.

    \item \textbf{\texttt{display\_lines.py}:} Cleans raw IR text into per-function line lists used by the rich animation panels.
\end{itemize}

\subsubsection{Transformation Layer}
The Transformation layer converts the ingest output into an intermediate scene graph and execution traces suitable for rendering.

\begin{itemize}
    \item \textbf{\texttt{models.py}:} Defines the core data model---\texttt{SceneGraph}, \texttt{SceneNode}, \texttt{SceneEdge}, \texttt{CFGBlock}, \texttt{TraceOverlay}, and \texttt{BlockMetadata}. Nodes carry a \texttt{role} assigned from edge topology: \texttt{entry}, \texttt{linear}, \texttt{branch}, \texttt{merge}, or \texttt{exit}.

    \item \textbf{\texttt{scene.py}:} \texttt{build\_scene\_graph(stream, ...)} groups events by function and block, attaches typed CFG edges, assigns block roles, applies optional domtree/loop metadata, and attaches optional trace overlays. Output is a complete \texttt{SceneGraph}.

    \item \textbf{\texttt{trace.py}:} \texttt{build\_execution\_trace(stream, include\_ssa=...)} produces a list of \texttt{TraceStep} or \texttt{RichTraceStep} records capturing the execution order with optional SSA result values for \texttt{binop}, \texttt{compare}, and \texttt{load} instructions.

    \item \textbf{\texttt{commands.py}:} \texttt{build\_animation\_commands(stream)} generates a list of \texttt{AnimationCommand} records (\texttt{push\_frame}, \texttt{pop\_frame}, \texttt{highlight\_cell}, etc.) used by the basic animation mode.
\end{itemize}

\subsubsection{Presentation Layer}
The Presentation layer renders the scene graph using Manim Community Edition\cite{manimce} and provides export functionality.

\begin{itemize}
    \item \textbf{User Interface:} The CLI accepts LLVM IR files and supports the following flags:
    \begin{lstlisting}[language=bash,caption={Actual CLI usage},gobble=4]
    # Export JSON scene graph + Graphviz DOT/PNG
    $ uv run llvmanim input.ll --json --draw \
        --outdir out/
    # Render basic stack animation
    $ uv run llvmanim input.ll --animate \
        --ir-mode basic --outdir out/
    # Render rich IR spotlight with SSA panel
    $ uv run llvmanim input.ll --animate \
        --ir-mode rich-ssa --speed 1.5 --outdir out/
    # CFG traversal animation
    $ uv run llvmanim input.ll --cfg-animate \
        --dot-cfg main.dot --import-trace trace.json \
        --outdir out/
    \end{lstlisting}

    When the tool encounters unsupported IR constructs, missing debug metadata, or malformed input, it emits a descriptive error message and continues processing the remainder of the program where possible.

    \item \textbf{Animation Modes:}
    \begin{itemize}
        \item \texttt{--ir-mode basic}: Stack-only view with per-cell badge flash for each instruction.
        \item \texttt{--ir-mode rich}: Two-column layout with an IR source panel and a moving spotlight cursor alongside the stack view.
        \item \texttt{--ir-mode rich-ssa}: Three-column layout (IR Source $\mid$ SSA Values $\mid$ Stack) showing computed results for binary operations, comparisons, and loads.
    \end{itemize}

    \item \textbf{CFG Traversal:} \texttt{--cfg-animate} renders an animated walk through the control-flow graph using DOT-derived node positions (\texttt{--dot-cfg}) and a runtime trace (\texttt{--import-trace}).

    \item \textbf{Export:}
    \begin{itemize}
        \item \texttt{--json}: Writes a machine-readable \texttt{scene\_graph.json}.
        \item \texttt{--draw}: Writes \texttt{cfg\_main.dot} and, when Graphviz binaries are available, \texttt{cfg\_main.png}.
    \end{itemize}

    \item \textbf{Rendering:} Manim CE renders MP4 or GIF output. GIF conversion uses an \texttt{ffmpeg} palette workflow to reduce peak memory usage. The \texttt{--preview} flag opens the rendered animation in a viewer.
\end{itemize}
